\section{Introduction}
In order to investigate the effects of interface type in touch-based gaming, a series of user experiments were conducted. This was necessary as the hypothesis under question examines a qualitative topic. That is, in order to determine which interface type is preferred, one needs to consider aspects related to enjoyment and gaming experience. 

Two weeks of user testing were conducted, with each team member responsible for conducting fifteen sessions with their game. The design of these experiments is described in more detail below. In addition the results are presented in Section 5.3, with a discussion of their findings provided in Section 5.4.


\section{Method}

\subsection{Independent and dependant variables}
In order to investigate the effect of gestures on touch-based gaming, three independent variables were used. These include interface type, game genre and participant characteristics. Interface type consists of three levels, namely direct manipulation, gestural and mixed interface versions. Game genre is composed of Puzzle, Strategy and Duelling/Fighting conditions. Finally, the participant characteristics that were investigated are the effect of gaming experience and previous use of a touch-device. These will be contrasted with the results of a control group who have neither significant gaming, nor touch-device experience.

In all cases, the dependant variable was the level of flow experienced by each participant. While one could also assess presence or immersion in the context of gaming, it was felt that flow would be most appropriate because the level of sensory stimulation might not be sufficient for presence to be experienced and no well-validated measures of immersion exist. In addition, as immersion is a necessary by-product of flow, a measure of flow should encapsulate this. Finally, as the experiment aims to assess which interface version elicits the most enjoyment, and flow is said to describe a state that is ``so gratifying that people are willing to do it for its own sake with little concern for what they will get out of it"~\citep{Csik90}, flow appears to be the most appropriate measure of gaming experience in this case.

\subsection{Experimental design}
The experiment consisted of a 3 (Interface type) $\times$ 3 (Genre) $\times$ 3 (Previous experience) repeated measures, two-way between subjects design. Interface type was assessed through repeated measures. That is, each participant played three versions of a game, each with a different interface type (direct manipulation, gestural and mixed). Participants were split into three mutually exclusive groups, where those in each group played a game of a different genre. Within each of these genre groups, participants were further subdivided according to various subject characteristics. These included experience with touch-devices and gaming experience. A third group was formed as a control for these subject characteristics.

\subsection{Instruments}
Two primary instruments were used in the execution of this experiment. These include a pre-experiment questionnaire used to assign participants to experimental conditions as well as a measure of flow.

\subsubsection{Pre-experiment questionnaire}
The pre-experiment questionnaire was used to assess which game genre and participant characteristic group a subject should be placed in. In addition, it was used to determine the level of computer literacy possessed by participants. The questionnaire has been included as Appendix~\ref{sec:preEx}. It consisted of a section gathering demographic information such as age, gender, the degree towards which the participant was studying and year of study. A 7-point Likert scale was used to assess computer literacy. 

In addition to demographic information, data was collected on both the participant characteristics under study as well as gaming preferences. This was done using a number of 7-point Likert scales. In order to assess the participant characteristics, questions regarding how often subjects played games as well as the frequency with which they interacted with touch-devices and how comfortable they felt using them, were asked.

Finally, the frequency with which participants played games within each genre as well as how much they enjoyed each genre was assessed. The scores for each of these sections were used to assign subjects to experimental groups. The manner in which this was done is discussed in Section 5.2.5.

\subsubsection{Flow assessment}
The Flow State Scales, developed by~\citet{Jack01}, were used to assess flow. This questionnaire uses nine factors in order to assess an individual's flow state. It was chosen due to its relatively short administration time. In addition, the instrument has been shown to have high reliability ($\alpha$ = 0.83)~\citep{Jack01} and internal consistency (Cronbach's $\alpha$ = 0.93)~\citep{Tene99}. Tests have also been done confirming the discriminant and construct validity of the measure~\citep{Mars99}.

While multiple versions of the test exist, the choice was made to use the \emph{Long General Flow State Scale}. The long scale was chosen as it provides a more thorough analysis of each factor. In addition, some research sheds doubt on the validity and reliability of the short scale~\citep{Tene99}. As the questionnaire was initially developed to assess flow in physical settings, the \emph{General} version was used as this is not specific to physical activity. Finally, the \emph{Flow State} rather than \emph{Dispositional} scale was used as the \emph{Dispositional} scale requires the test to be administered during the execution of some flow-inducing activity. Due to the nature of the experiment, there were concerns over the time taken to enter a flow state. In addition, it was felt that answering a questionnaire during game-play would disrupt any flow state achieved. Thus the choice was made to administer the questionnaire after each game had been completed.

\begin{table}[!ht]

\centering %
\begin{tabular*}{\textwidth}{@{\extracolsep{\fill}} p{4cm} p{11cm} }
	\hline	
	 Flow dimension & Description \\
	\hline
	\emph{Challenge-Skill Balance} & This describes an appropriate balance between the difficulty presented by a task and the ability or perceived ability of an individual to complete this task. \\ 
	\emph{Merging of Action and Awareness} & This refers to a state of deep immersion where an individual feels at one with their actions. \\
	\emph{Clear Goals} & A clarity of intention, aided by goals being set prior to beginning an activity. \\
	\emph{Unambiguous Feedback} & This refers to feedback that provides a clear awareness of how well one is progressing towards one's goals. \\
	\emph{Concentration on the Task at Hand} & A state where one has a single-minded focus on the task at hand, while ignoring any potential distractions. \\
	\emph{Sense of Control} & A feeling of total mastery over one's environment and actions. \\
	\emph{Loss of Self-Consciousness} & Concern for the self and other everyday issues disappears due to total concentration on the task at hand. \\
	\emph{Transformation of Time} & This describes a perceived distortion in the normal passage of time. \\
	\emph{Autotelic Experience} & This describes an experience that is intrinsically rewarding and not simply a means to an end. \\
	\hline
	\end{tabular*}
\caption{Summary of the nine flow dimensions measured by the Flow State Scales developed by \citet{Jack01}.}
\label {flowDimen}
\end{table}

The questionnaire provided flow scores along nine dimensions, summarised in Table~\ref{flowDimen}. In addition, these factors can be combined to create an overall flow state score. Thus the questionnaire provides nine dependant measures of flow with an additional higher-order factor. The questionnaire consisted of 36 questions that assessed the various aspects of flow using a 5-point Likert scale, ranging from \emph{Strongly disagree} to \emph{Strongly agree}. These questions asked participants to assess various aspects of their mental state during the task. A select portion of the questionnaire has been included as Appendix~\ref{sec:flowScale}.



\subsection{Materials and setting}

Each participant played three versions of the same game. Three games were developed for this purpose. These games consisted of a Puzzle, Strategy and Fighting/Duelling game. Within each game genre, three versions of the same game were developed in order to test the effects of different interaction techniques (direct manipulation, gestures and a combination of the two). Thus, nine game versions were used in total. However, subjects only played three of these games each, with all three belonging to the same game genre. Participants played the games on the Apple iPad. This was chosen due to its multi-touch capabilities as well as large screen size when compared to other commercially available multi-touch devices. While the Microsoft Surface has a larger screen, its prohibitive cost means that this device would not commonly be available to gamers which would reduce the ecological validity of our study to some degree.

\subsubsection{Experiment setting}
All participants completed the experiment in the same room in the Mezzanine level of the Computer Science building. The room contained only a table and the iPad, in order to minimise distractions. In addition, noise levels were kept consistently low during experimentation. Participants were placed in this room alone, with an experimenter available outside the room at all times should the participant have any questions regarding the task at hand. 

\subsection{Participants}
Non-probabilistic, quota sampling was used to recruit 45 participants. Each participant was paid R20 for completing the experiment. Participants were also required to sign a consent form which explained their rights with regard to the study and also confirmed their understanding of what was required of them. This consent form is provided as Appendix~\ref{sec:consent}. All of the recruited participants were University of Cape Town students, the majority of whom were studying towards a degree in Computer Science, Engineering or a similar Science related field. Their average year of study was 2.4 (\emph{SD} = 1.48). Of the participants, 35 were male and 10 female, with their age ranging from 18 to 29 (\emph{M} = 20.98, \emph{SD} = 2.01).

Of the 45 participants, 15 were assigned to each game genre. Of these 15, 5 were matched to each of the two participant characteristic conditions, with the remaining 5 serving as the control group. The number of participants was chosen based on recommendations that a ratio of at least five participants be observed for each independent variable~\citep{Bart01}. The choice of how to assign participants to each group was determined through the use of the pre-experiment questionnaire. This questionnaire assessed participants' levels of gaming and touch-device experience. Table~\ref{allGameTouch} summarises the scores for these questions within each participant characteristic group. In addition, it also provided a measure of what game genre participants enjoyed most of those being tested. This was done so that participants were assigned to a game genre they had expressed interest in, in order to prevent a simple dislike of the type of game being played from affecting results. 

\begin{table}
\centering %
\begin{tabular*}{0.75\textwidth}{@{\extracolsep{\fill}} l p{1.25cm} p{1.25cm} p{1.25cm} p{1.25cm}  }

\hline	
  & \multicolumn{2}{p{4cm}}{\emph{Touch-device experience}} & \multicolumn{2}{p{4cm}}{\emph{Gaming experience}}   \\
  & \emph{M} & \emph{SD} & \emph{M} & \emph{SD}  \\
  \hline
 Control group & 3.2 & 1.26 & 3.8 & 1.24 \\
 Touch experience group & 5.27 & 0.51 & 4.53 & 1.45 \\
 Gaming group & 4 & 1.49 & 6.07 & 0.7 \\
\hline  
\end{tabular*}%}
\caption{Descriptive statistics of participant characteristics within each experience group across genre. All results were obtained using a 7-point Likert scale.}
\label {allGameTouch}
\end{table}

An attempt was made to control for computer literacy amongst participants as differences in this variable could foreseeably influence the results achieved. This was done by recruiting participants from computer laboratories in the Computer Science building. In addition, a question was included on the pre-experiment questionnaire in order to assess this factor. Participants who did not score above the midpoint (4) on this measure were disregarded for the study. On the 7-point scale, the mean computer literacy score for all participants was 6.2 (\emph{SD} = 0.77). No attempt was made to balance for factors such as race, age or gender as these were not felt to be particularly relevant to the topic at hand. In addition, controlling for these elements is beyond the scope of this study as it would have required a far greater number of participants.

\subsubsection{Participant characteristics of Puzzle group}
The Puzzle game group consisted of 15 participants. Of these, nine were male and six female, ranging in age from 18 to 24 (\emph{M} = 20.8, \emph{SD} = 1.74). Nine of these subjects were completing Computer Science related degrees with the remaining participants studying other Science disciplines or Accounting. The year of study ranged between 1 and 4 (\emph{M} = 2.4, \emph{SD} = 0.99). All individuals rated themselves as being highly computer literate on the 7-point scale (\emph{M}=6.2, \emph{SD} = 0.77). Table~\ref{touchGaming} summarises the scores on the two participant variables within the gaming experience, touch-device experience and control groups. Participants were assigned to this group due to their expressed interest in puzzle games. The mean score when subjects were asked how much they enjoyed puzzle games was 5.8 (\emph{SD} = 1.4).

\begin{table}
\centering %
\begin{tabular*}{0.75\textwidth}{@{\extracolsep{\fill}} l p{1.25cm} p{1.25cm} p{1.25cm} p{1.25cm}  }

\hline	
  & \multicolumn{2}{p{4cm}}{\emph{Touch-device experience}} & \multicolumn{2}{p{4cm}}{\emph{Gaming experience}}   \\
  & \emph{M} & \emph{SD} & \emph{M} & \emph{SD}  \\
  \hline
 Control group & 2.4 & 1.14 & 3.2 & 1.1 \\
 Touch experience group & 5.8 & 0.45 & 4 & 2.12 \\
 Gaming group & 3.8 & 1.64 & 5.8 & 1.1 \\
\hline  
\end{tabular*}
\caption{Descriptive statistics of participant characteristics within the three groups that compose the Puzzle game condition. All results were obtained using a 7-point Likert scale.}
\label {touchGaming}
\end{table}

\subsection{Control procedures}
Repeated measures was used to assess interface type as it removes the need for matching between groups. However, using a repeated measures design introduces a number of caveats that must be handled. These can be classed as order effects and include fatigue, contrast and practice effects. The literature suggests using either counterbalancing or allowing enough time between each condition so that the order effects are reduced~\citep{Cozb05}. An attempt was made to use both techniques. That is, the order with which participants complete the experimental conditions was varied through randomisation. In addition, participants were required to complete the Flow State questionnaire between each condition. This increased the amount of time between each condition. A distractor task was not used as it would have increased the overall length of the experiment dramatically, perhaps introducing fatigue effects.

A tutorial session was included so that the effect of unfamiliarity with either gestures, or touch-devices in general, could be minimised. In this session, participants were instructed to draw the gestures that would be used in the game a certain number of times, with the tutorial only progressing once the gesture had been successfully recognised. This served to familiarise participants with both the system and this style of interaction. While this session may have reduced between-group differences with regard to touch-device experience, it was felt to be necessary as difficulty in using or understanding how the gestural interface worked would have severely compromised the degree to which flow was elicited in this condition. In addition, the practice session should not have proved sufficient to provide participants with enough experience with touch-devices for this participant characteristic to have been strongly influenced.

The experimental system was designed so as to minimise experimenter bias. Instructions on how to complete the tasks were given to participants in written form and included in the games. This was done to decrease the degree of experimenter-participant interaction. More importantly, the randomisation and execution of each game was entirely automated. That is, a system was written to link all three games, the tutorial session and all instructions together so that no experimenter interaction was required throughout the experiment. As the experiment does not lend itself to single- or double-blind conditions, since participants are required to be aware of which condition they are completing in order to play the game, the automated system should remove the majority of experimenter bias. Attempts were also made to keep experimenter behaviour consistent before the experiment began. Finally, in order to reduce the number of demand characteristics, the purpose of the study was not explained to participants until they had completed it. This was done as any prior knowledge of the hypothesis under question may have influenced participants to attempt to either confirm or disconfirm it.

Finally, a number of situational effects were also considered. Participants all completed the study in the same setting with minimal distractions or external elements. In addition, in order to prevent the Hawthorne effect, participants were not observed during the experiment. That is, all participants played the game in a separate room. The experimenter was present outside this room should they have had any questions. The choice to not watch participants (rather than to watch all participants) was made  as observation may decrease subject performance due to nervousness and other such factors.

\subsection{Procedure}
The following section describes the experimental procedure followed for each participant. Participants were seated and asked to sign a consent form and complete a demographic information sheet. They were then given written instruction on how to complete the experiment. These forms can be found in Appendix~\ref{sec:consent}, \ref{demographic} and \ref{introInstruct}, respectively.

The experiment consisted of seven steps. Initially participants completed a tutorial session where they were asked to practice drawing the gestures that would be used in the game. After this they were instructed on how to play the first interface version of the game. Game-play lasted seven minutes. This time frame was controlled by a timer within the game. After this time had elapsed, game-play was interrupted and the participant was instructed to complete the first questionnaire. This process of instructions, seven minutes of game-play and then filling out the questionnaire was repeated for each of the other two game versions. Once all three games had been played, the participant was informed that the experiment was over and thanked for their participation. The entire process was automated and required no experimenter interaction. Once the participant had completed the experiment, they were paid R20 and given a sheet explaining the purpose of the experiment. This form is available in Appendix~\ref{postEForm}. They were also able to ask any questions they may have had regarding the topic under study. The entire procedure took roughly 45 minutes to complete.

\subsection{Pilot study}
A pilot study was held in order to assess the overall experimental procedure. Five participants were recruited from the Computer Science Honours class. Each participant played through the combined system which linked together all three game versions with the tutorial screen and timing and logging code. The participants also completed the questionnaire between each game version. This was done so that an approximation of total experiment time could be made. Initially each game was played for ten minutes. However, upon reviewing total experiment time, the decision was made to reduce this to seven minutes per game, as the experiment was initially felt to be too lengthy. The pilot study also allowed the logging and randomisation code to be checked. Finally, participants were asked to provide comments on the thoroughness of the in-game instructions and tutorial text, as well as any general comments they may have had on the experimental procedure. Following this, small adjustments were made to the tutorial screens to make them more informative. The overall design of the experiment remained unchanged, however.
